% ============================================================
% 封面、前言、广东卷命题分析
% ============================================================
\begin{titlepage}
\begin{tikzpicture}[remember picture,overlay]
\fill[mainblue] (current page.north west) rectangle ([yshift=-9.5cm]current page.north east);
\fill[mainblue!80!black] ([yshift=-9.5cm]current page.north west) rectangle ([yshift=-10.1cm]current page.north east);
% 装饰：苯环与分子
\begin{scope}[shift={([xshift=-4.2cm,yshift=-4.6cm]current page.north east)},white,opacity=0.35,line width=1.2pt]
\foreach \a in {0,60,...,300} {\draw (\a:1.6)--(\a+60:1.6);}
\draw (0,0) circle (1.0);
\foreach \a in {30,150,270} {\draw (\a:1.6*0.866*1.15)--(\a:2.6);}
\end{scope}
\begin{scope}[shift={([xshift=3.2cm,yshift=-7.6cm]current page.north west)},white,opacity=0.25,line width=1pt]
\foreach \x/\y in {0/0,1.2/0.7,2.4/0,3.6/0.7,4.8/0}{\fill (\x,\y) circle (0.18);}
\draw (0,0)--(1.2,0.7)--(2.4,0)--(3.6,0.7)--(4.8,0);
\end{scope}
\node[anchor=west,text=white,font=\heiti\fontsize{40}{48}\selectfont] at ([xshift=2.2cm,yshift=-3.9cm]current page.north west) {高中化学系统讲义};
\node[anchor=west,text=white,font=\kaishu\Large] at ([xshift=2.3cm,yshift=-5.5cm]current page.north west) {选择性必修 1·2·3 与高考综合};
\node[anchor=west,text=white!85!mainblue,font=\heiti\large] at ([xshift=2.3cm,yshift=-6.6cm]current page.north west) {面向广东新高考\ ·\ 目标 90 分};
\end{tikzpicture}

\vspace*{10.2cm}
\begin{center}
\begin{minipage}{0.86\linewidth}
\begin{tcolorbox}[enhanced,colback=lightblue,colframe=mainblue,boxrule=0.8pt,arc=2mm,left=4mm,right=4mm,top=3mm,bottom=3mm]
\heiti\large 全书结构\normalfont\normalsize\par\medskip
\begin{tabular}{@{}ll@{}}
\textbf{预备篇} & 第0章\quad 高一必备基础回顾 \\[2pt]
\textbf{原理篇} & 第1章 热效应\quad 第2章 速率与历程\quad 第3章 化学平衡 \\
& 第4章 水溶液中的离子平衡\quad 第5章 化学反应与电能 \\[2pt]
\textbf{结构篇} & 第6章 原子结构\quad 第7章 分子结构\quad 第8章 晶体结构 \\[2pt]
\textbf{有机篇} & 第9章 结构与烃\quad 第10章 烃的衍生物\quad 第11章 大分子与合成 \\[2pt]
\textbf{综合篇} & 第12章 化学实验综合\quad 第13章 元素化合物与化工流程 \\
& 第14章 高考综合大题训练 \\
\end{tabular}
\end{tcolorbox}
\vspace{1.2cm}
{\kaishu 每章由正文讲解、例题、本章小结和分层练习组成，练习附有完整解析。}
\end{minipage}
\end{center}
\vfill
\begin{center}
{\heiti\large 2026 年 10 月版}
\end{center}
\end{titlepage}

\chapter*{前\quad 言}
\markboth{前言}{}
\addcontentsline{toc}{chapter}{前言}

这本讲义写给学完高一化学（必修第一、二册）、准备学习选择性必修内容并参加广东新高考的同学。书中假定你已经学过物质的量、离子反应、氧化还原反应和元素周期律，但还没有学过高二的任何内容；凡是用到高一知识的地方，书中会先简短回顾，再往下讲。

全书分为五篇。预备篇只有第0章，复习后面各章离不开的高一基础。原理篇（第1～5章）对应选择性必修1，依次讲化学反应的热效应、反应速率、化学平衡、水溶液中的离子平衡和电化学，这五章前后衔接紧密，后一章要用到前一章的概念和方法。结构篇（第6～8章）对应选择性必修2，讲原子、分子和晶体的结构；有机篇（第9～11章）对应选择性必修3。综合篇（第12～14章）讨论化学实验、元素化合物与工艺流程，并给出四道按广东卷第17～20题题型编写的综合大题，适合在三本选择性必修都学完之后使用。

每章开头有一段引言，说明本章研究什么、和前后各章有什么联系。正文讲一个新概念时，一般先说明研究对象，再从熟悉的现象或图示入手建立直观认识，然后给出规范的定义或规律，最后配上由浅入深的例题。高中阶段能够讲清楚的道理，书中尽量讲清楚；有些问题要用到大学的理论才能深入解释（例如电子运动的量子力学描述、熵的统计意义），书中只说明高中采用的是什么模型、在什么范围内适用。章末有本章小结和分层练习：A 组检查基本概念和基本计算，B 组是高考常见难度，C 组综合性较强，所有练习都附有解析。

建议的学习顺序是：先读第0章，并用章末 A 组练习检查自己的高一基础，做错的题回到对应小节重读；然后按顺序学习原理篇。结构篇和有机篇可以按学校的教学进度安排，二者相对独立，但有机篇中碳原子的杂化方式、分子的极性和氢键要用到第6、7章，最好先学这两章。综合篇放在最后。

读例题时，先遮住解析自己做，做完再对照。对照时不要只看答案是否一致，还要看解析中每一步的依据：为什么选这个公式，为什么可以忽略某个量，为什么这样判断平衡移动的方向。练习也要先独立完成再看解析；做错的题记下错因，过一段时间再做一遍。

书中标“真题”的题目取自 2024、2025、2026 年广东高考化学试卷，并注明年份和题号；其中 2026 年试卷依据的是网络收集版，个别文字可能与正式版本有出入。对真题题干作了删改的，标为“改编”；其余题目为原创，数据都经过验算。真题中的图片在书中用文字或重新绘制的示意图转述。

\chapter*{广东卷命题分析}
\markboth{广东卷命题分析}{}
\addcontentsline{toc}{chapter}{广东卷命题分析}

\section*{一、试卷结构}
广东新高考化学满分 100 分，考试时间 75 分钟。第Ⅰ卷为 16 道单项选择题，共 44 分（第 1～10 题每题 2 分，第 11～16 题每题 4 分）；第Ⅱ卷为 4 道非选择题，每题 14 分，共 56 分，题型相对固定：第17题化学实验探究，第18题工艺流程，第19题化学反应原理综合，第20题有机化学综合。

\section*{二、2024—2026 年广东卷逐题考点}
下表按题号列出三年试卷的情境和主要考点。横向看同一题号，可以看出哪些题位的考查内容比较固定。
\begin{center}
\small
\renewcommand{\arraystretch}{1.12}
\begin{tabularx}{\linewidth}{c>{\raggedright\arraybackslash}X>{\raggedright\arraybackslash}X>{\raggedright\arraybackslash}X}
\toprule
题号 & 2024 年 & 2025 年 & 2026 年（网络收集版） \\
\midrule
1 & 文物材料（有机高分子） & 传统工艺材料（合金） & 桥梁构件材质（有机物） \\
2 & 科技成就：变化、能量、电子式 & 多烯环配合物：共价键、大 $\pi$ 键 & 科技成就：放热、核素、周期 \\
3 & 计时器材料：合金、石英、核素 & 科技情境：摩尔质量、化合价、能量转化 & 竹：物理变化、漂白、吸附 \\
4 & 饮食：糖类、油脂、极性分子 & 劳动项目：脱硫、盐析、铁的氧化 & 成像：能量转化、\ce{SiO2}、共价晶体 \\
5 & 金属防腐（牺牲阳极法） & 存储材料：单体、\ce{Fe3O4}、晶体 & 陈述关联：胶体、晶体、焰色 \\
6 & 苯甲酸重结晶 & 铁钉镀铜实验操作 & 海上风电防腐 \\
7 & 劳动项目与化学知识 & 物质转化链（S、Na、C、Al） & 阿伏加德罗常数 \\
8 & \ce{Cl2} 的制备与性质装置 & 陈述因果（浓硫酸、\ce{Fe^2+} 检验） & 结构解释性质（氢键、推电子效应） \\
9 & 陌生有机物结构与性质 & 第一电离能图与周期律 & \ce{Cl2} 制备及与 Fe 反应的实验 \\
10 & 阿伏加德罗常数 & 阿伏加德罗常数 & 岭南美食中的化学 \\
11 & 吸热反应、熵增、\ce{NH3} 实验 & \ce{H2O2} 分解与乙醇催化氧化实验 & 元素推断（3p 半充满） \\
12 & 元素推断（结构式） & CuCl 配位平衡、电荷守恒 & 卤代烃消去与水解实验 \\
13 & 陈述因果（酸性、冠醚） & 由结构推测性质 & 价类二维图（C、N、Cl） \\
14 & 价类二维图（Mg、Al、Fe） & 水系可充电电池 & 方程式正误 \\
15 & 反应历程能量图（中间体、决速步） & Cu 与稀硝酸反应实验 & 平行反应能量图与 $c$–$t$ 图 \\
16 & 离子交换膜电解与计算 & 燃料电池催化机理 & 电解还原与海水脱盐 \\
17 & 硫的转化；滴定；速率与平衡建立 & 重结晶；pH 滴定求 $K_\text{a}$；方案设计 & 氨水制备；镀铜；铜氨配合物的 pH 滴定 \\
18 & 炭渣提镓：配离子、$K_\text{a}$ 与 $K_\text{sp}$、晶胞 & 多金属精矿：$K_\text{sp}$、萃取、晶胞、电解 & 金精矿提金：电解、$K_\text{sp}$、配平、超分子、晶胞 \\
19 & 轨道表示式；盖斯定律；催化；非水溶剂中的 $K_\text{a}$ & 价层电子排布；$\Delta G$；盖斯定律；平衡组成与 $K$；选择性 & 轨道表示式；盖斯定律；速率数据；熵变；吸收量与 $K$ \\
20 & 苯乙烯$\to$手性醇；同分异构；合成 & 光酶三组分反应；核磁共振氢谱；合成 & 羟基迁移；原子利用率；合成路线 \\
\bottomrule
\end{tabularx}
\end{center}

\section*{三、高频考点与命题趋势}
\begin{enumerate}
\item 有些考点每年都出现：阿伏加德罗常数（2024、2025 年第10题，2026 年第7题），元素推断与周期律，价类二维图或物质转化，两道电化学选择题（防腐、电池或电解）。反应历程能量图在 2024、2026 年的第15题出现。四道大题的题型基本固定。这些内容是复习时首先要保证得分的部分。
\item 第19题（原理综合）三年都考了盖斯定律求 $\Delta H$、平衡常数或吸收量的计算，也都有一个结构小问（轨道表示式或价层电子排布式），并且都要求分析速率或选择性数据。计算题常要求“写出计算过程”，书写要规范。
\item 第18题（工艺流程）综合性强，三年都考了 $K_\text{sp}$ 计算、晶胞、陌生方程式书写和配合物（或超分子），一道题同时涉及元素化合物、反应原理和物质结构三方面的知识。
\item 第17题（实验探究）偏重定量和探究：三年都有滴定计算（2025、2026 年是 pH 滴定），并按“提出假设—设计验证—得出结论”的过程设问。
\item 第20题（有机综合）以新的合成方法为情境（3D 打印催化材料、光酶催化、羟基迁移），考查官能团、同分异构体、杂化与手性、合成路线设计，要求把题中的新反应迁移到新的底物上。
\item 选择题前几题多以传统文化、科技成就、劳动项目、岭南生活为情境；每年都有“陈述Ⅰ与陈述Ⅱ”或“结构解释性质”式的因果判断题。
\end{enumerate}

\section*{四、各题型的得分目标}
以 90 分为目标，各部分可以按下表分配。表中最后一列是这部分题目的主要难点和本书中对应的章节。
\begin{center}
\small
\begin{tabularx}{\linewidth}{p{3.2cm}cX}
\toprule
板块 & 目标得分 & 主要难点 \\
\midrule
选择题 1～10（20 分） & 20 & 概念准确：STSE、$N_\text{A}$、实验操作、结构与性质（第0、6～8、12章） \\
选择题 11～16（24 分） & $\ge20$ & 图像与电化学：能量图、速率与平衡图像、电池与电解池（第2～5章） \\
第17题 实验（14 分） & $\ge12$ & 操作规范、滴定计算、假设与方案的表述（第12章） \\
第18题 流程（14 分） & $\ge12$ & 预处理与条件控制、方程式配平、$K_\text{sp}$、晶胞（第4、8、13章） \\
第19题 原理（14 分） & $\ge12$ & 盖斯定律、三段式、$K_p$、多重平衡与选择性（第1、3章） \\
第20题 有机（14 分） & $\ge12$ & 同分异构体计数、合成路线设计（第9～11章） \\
\bottomrule
\end{tabularx}
\end{center}
